\section*{Bab \ref{ch:b2:enolates-aldol} --- Enol, Enolat, dan Reaksi Aldol}

\begin{solution}{exo:b2:enolates-aldol:1}
Butanon: but-1-en-2-ol \ce{CH2=C(OH)CH2CH3} dan but-2-en-2-ol \ce{CH3C(OH)=CHCH3} ((\textit{E}) dan
(\textit{Z})). Sikloheksanon: sikloheks-1-en-1-ol.
\end{solution}

\begin{solution}{exo:b2:enolates-aldol:2}
Etana $<$ propanon $<$ etil 3-oksobutanoat (10.7) $<$ pentana-2,4-dion (9.0).
\end{solution}

\begin{solution}{exo:b2:enolates-aldol:3}
3-Hidroksibutanal, lalu but-2-enal pada pemanasan.
\end{solution}

\begin{solution}{exo:b2:enolates-aldol:4}
Dietil propanadioat, natrium etoksida, 1-bromobutana; hidrolisis menjadi
asam butilpropanadioat; pemanasan melepaskan \ce{CO2}:
\ce{CH3(CH2)3CH2COOH}, asam heksanoat.
\end{solution}

\begin{solution}{exo:b2:enolates-aldol:5}
Kinetik: deprotonasi pada C6 (sisi \ce{CH2}), dengan LDA pada \qty{-78}{\degreeCelsius}. Termodinamik:
ikatan rangkap C1=C2 yang lebih tersubstitusi, ke arah karbon yang membawa
metil, dengan alkoksida dalam alkoholnya pada suhu kamar.
\end{solution}

\begin{solution}{exo:b2:enolates-aldol:6}
4-Fenilbut-3-en-2-on (benzalaseton), \ce{C6H5CH=CHCOCH3}. Benzaldehida tidak
mempunyai hidrogen $\alpha$: senyawa ini tidak dapat membentuk \omterm{def:b2:enolates-aldol:enolate}{enolat} dan
hanya menjadi elektrofil.
\end{solution}

\begin{solution}{exo:b2:enolates-aldol:7}
\omterm{def:b2:enolates-aldol:enolate}{Enolat} salah satu gugus metil (C1) menyerang karbonil yang lain (C5): cincin
beranggota lima, 3-metilsiklopent-2-en-1-on setelah dehidrasi. Alternatifnya,
\omterm{def:b2:enolates-aldol:enolate}{enolat} pada C3 yang menyerang C5, akan membentuk cincin beranggota tiga yang
tegang.
\end{solution}

\begin{solution}{exo:b2:enolates-aldol:8}
\ce{2CH3COOC2H5 -> CH3COCH2COOC2H5 + C2H5OH} dengan natrium etoksida; produknya,
yang terdeprotonasi, diperoleh kembali sebagai etil 3-oksobutanoat pada
pengasaman.
\end{solution}

\begin{solution}{exo:b2:enolates-aldol:9}
Propanon, etanal, dan asetofenon (ketiganya mempunyai gugus \ce{CH3-CO};
etanal adalah \ce{CH3CHO}). Pentan-3-on tidak.
\end{solution}

\begin{solution}{exo:b2:enolates-aldol:10}
\omterm{def:b2:enolates-aldol:enolate}{Enolat} di salah satu ujung menyerang ester di ujung yang lain: etil
2-oksosiklopentana-1-karboksilat, cincin beranggota lima.
\end{solution}

\begin{solution}{exo:b2:enolates-aldol:11}
Dalam 3-metilpentan-2-on, pusat stereogeniknya (C3) adalah karbon $\alpha$:
enolisasi membuatnya planar, dan reprotonasi dari salah satu muka
merasemisasinya. Dalam 4-metilheksan-2-on, pusat stereogeniknya berada pada
C4, $\beta$ terhadap karbonil, dan tidak tersentuh oleh enolisasi.
\end{solution}

\begin{solution}{exo:b2:enolates-aldol:12}
Propanon dengan dua ekuivalen benzaldehida dalam basa. Dengan dua aldehida
yang berbeda, kondensasikan propanon dengan salah satu aldehida lebih
dahulu (satu ekuivalen, dengan mengisolasi benzalaseton), lalu
kondensasikan gugus metilnya yang tersisa dengan aldehida kedua. Dalam satu
wadah, kedua aldehida itu bersaing dan menghasilkan campuran tiga dienon.
\end{solution}

\begin{solution}{pb:b2:enolates-aldol:1}
\textbf{1.} Propanon, enam hidrogen $\alpha$ (tiga pada setiap metil).
\textbf{2.} $\mathrm{CH_3{-}CO{-}CH_2^-} \leftrightarrow \mathrm{CH_3{-}C(O^-){=}CH_2}$.
\textbf{3.} Benzaldehida tidak mempunyai hidrogen pada karbon di sebelah
karbonilnya (karbon itu termasuk dalam cincin dan tidak membawa H).
\textbf{4.} Benzaldehida, sebuah aldehida yang berlebih, adalah elektrofil
yang lebih baik; \omterm{def:b2:enolates-aldol:enolate}{enolat} propanon lebih sering bertemu dengannya daripada
dengan karbonil propanon.
\textbf{5.} Keton mempunyai dua gugus karbon yang mendorong kerapatan
elektron ke karbon karbonilnya dan menghalangi pendekatan; aldehida hanya
satu.
\textbf{6.} Hanya sebagian kecil \omterm{def:b2:enolates-aldol:enolate}{enolat} yang diperlukan: \omterm{def:b2:enolates-aldol:enolate}{enolat} itu
diregenerasi sambil bereaksi.
\textbf{7.} Karbon \omterm{def:b2:enolates-aldol:enolate}{enolat} berikatan dengan karbon karbonil benzaldehida:
alkoksida 4-hidroksi-4-fenilbutan-2-on.
\textbf{8.} Pelepasan sebuah hidrogen $\alpha$ dan lepasnya hidroksida dari karbon $\beta$: 4-fenilbut-3-en-2-on.
\textbf{9.} Ikatan rangkap yang terbentuk terkonjugasi dengan cincin sekaligus dengan karbonil.
\textbf{10.} Gugus metil yang lain masih mempunyai tiga hidrogen $\alpha$.
\textbf{11.} \ce{2C6H5CHO + CH3COCH3 -> C6H5CH=CHCOCH=CHC6H5 + 2H2O}.
\textbf{12.} Dehidrasi-dehidrasinya, dan pengendapan produk, yang
meninggalkan larutan.
\textbf{13.} Dua, keduanya (\textit{E}).
\textbf{14.} Isomer-isomer (\textit{Z}) menempatkan gugus fenil di sebelah
gugus karbonil: tegangan sterik.
\textbf{15.} Dua \ce{C=C}, satu \ce{C=O}, dan kedua cincin: sistem
terkonjugasi yang membentang melintasi molekul.
\textbf{16.} Konjugasinya yang panjang mempersempit celah $\pi$ (\cref{prop:b2:huckel:gap}):
senyawa itu menyerap di ujung biru daerah tampak dan terlihat kuning.
\textbf{17.} Tidak: senyawa itu planar (atau hampir planar) dengan bidang cermin.
\textbf{18.} 106.12, 58.08, dan \qty{234.30}{g/mol}.
\textbf{19.} Benzaldehida $2.1 \times 1.045/106.12 = \qty{20.7}{mmol}$; propanon $0.75 \times
0.7845/58.08 = \qty{10.1}{mmol}$.
\textbf{20.} Propanon memerlukan \qty{20.3}{mmol} benzaldehida, dan yang
tersedia \qty{20.7}{mmol}: propanon menjadi pembatas, dengan selisih tipis.
\textbf{21.} $0.01013 \times 234.30 = \qty{2.37}{g}$.
\textbf{22.} Sisa benzaldehida tetap berada dalam etanol dan tercuci.
\textbf{23.} Produk monokondensasi, benzalaseton.
\textbf{24.} $0.75 \times 2.37 = \boldsymbol{\approx \qty{1.78}{g}}$ dibenzalaseton.
\end{solution}
